\documentclass{article}

\usepackage[12pt]{extsizes}
\usepackage[T2A]{fontenc}
\usepackage[utf8]{inputenc}
\usepackage[english, russian]{babel}

\usepackage{mathrsfs}
\usepackage[dvipsnames]{xcolor}

\usepackage{amsmath}
\usepackage{amssymb}
\usepackage{amsthm}
\usepackage{indentfirst}
\usepackage{amsfonts}
\usepackage{enumitem}
\usepackage{graphics}
\usepackage{tikz}
\usepackage{tabu}
\usepackage{diagbox}
\usepackage{hyperref}
\usepackage{mathtools}
\usepackage{ucs}
\usepackage{lipsum}
\usepackage{geometry} % Меняем поля страницы
\usepackage{fancyhdr} % Headers and footers
\newcommand{\range}{\mathrm{range}}
\newcommand{\dom}{\mathrm{dom}}
\newcommand{\N}{\mathbb{N}}
\newcommand{\R}{\mathbb{R}}
\newcommand{\E}{\mathbb{E}}
\newcommand{\D}{\mathbb{D}}
\newcommand{\M}{\mathcal{M}}
\newcommand{\Prime}{\mathbb{P}}
\newcommand{\A}{\mathbb{A}}
\newcommand{\Q}{\mathbb{Q}}
\newcommand{\Z}{\mathbb{Z}}
\newcommand{\F}{\mathbb{F}}
\newcommand{\CC}{\mathbb{C}}

\DeclarePairedDelimiter\abs{\lvert}{\rvert}
\DeclarePairedDelimiter\floor{\lfloor}{\rfloor}
\DeclarePairedDelimiter\ceil{\lceil}{\rceil}
\DeclarePairedDelimiter\lr{(}{)}
\DeclarePairedDelimiter\set{\{}{\}}
\DeclarePairedDelimiter\norm{\|}{\|}

\renewcommand{\labelenumi}{(\alph{enumi})}

\newcommand{\smallindent}{
    \geometry{left=1cm}% левое поле
    \geometry{right=1cm}% правое поле
    \geometry{top=1.5cm}% верхнее поле
    \geometry{bottom=1cm}% нижнее поле
}

\newcommand{\header}[3]{
    \pagestyle{fancy} % All pages have headers and footers
    \fancyhead{} % Blank out the default header
    \fancyfoot{} % Blank out the default footer
    \fancyhead[L]{#1}
    \fancyhead[C]{#2}
    \fancyhead[R]{#3}
}

\newcommand{\dividedinto}{
    \,\,\,\vdots\,\,\,
}

\newcommand{\littletaller}{\mathchoice{\vphantom{\big|}}{}{}{}}

\newcommand\restr[2]{{
    \left.\kern-\nulldelimiterspace % automatically resize the bar with \right
    #1 % the function
    \littletaller % pretend it's a little taller at normal size
    \right|_{#2} % this is the delimiter
}}

\DeclareGraphicsExtensions{.pdf,.png,.jpg}

\newenvironment{enumerate_boxed}[1][enumi]{\begin{enumerate}[label*=\protect\fbox{\arabic{#1}}]}{\end{enumerate}}



\smallindent

\header{ЦРОД $\bullet$ Математика}{\textit{Алгебра}}{2 ноября 2023}

\fancyhead[L]{Математика}
\fancyhead[C]{\textit{Алгебра}}
\fancyhead[R]{2 ноября 2023}% Custom header text

%----------------------------------------------------------------------------------------

\begin{document}
    \large

    \begin{center}
        \textbf{Интерполяция}
    \end{center}

    \epigraph{\textit{Повторение --- мать учения}}{}
    \begin{enumerate_boxed}
        \item Даны различные числа $x_0, \dotsc , x_n.$
        \begin{enumerate}
            \item Пусть многочлен $P_k(x)$ степени не больше $n$ равен $0$ при всех $x_0, \dotsc , x_n,$ кроме $x_k$.
            Докажите, что \[P_k(x) = c(x - x_0)\dotsc(x - x_{k-1})(x - x_{k+1})\dotsc(x - x_n).\]

            \item Чему должна быть равна константа $c$, чтобы $P_k(x_k)$ было равно $1$?

            \item Интерполяционный многочлен Лагранжа.
            Даны (не обязательно различные) числа $y_0, \dotsc , y_n$.
            Докажите, что единственный многочлен степени не выше $n$, принимающий в точках $x_0, \dotsc , x_n$ значения $y_0, \dotsc , y_n$, соответственно,
            равен

            \[\sum_{k = 0}^{n} y_k \left(\frac{x - x_0}{x_k - x_0}  \dotsc  \frac{x - x_{k-1}}{x_k - x_{k-1}} \frac{x - x_{k+1}}{x_k - x_{k+1}} \dotsc \frac{x - x_{n}}{x_k - x_{n}}\right) = y_{0}P_0(x) + \dotsc + y_{n}P_n(x)\]
        \end{enumerate}

        \item Найдите многочлен $f$ наименьшей степени такой, что \[f(-1) = 2, f(0) = 0, f(2) = 2, f(3) = 6.\]

        \item Найдите многочлен $f$ наименьшей степени такой, что \[f(0) = 0, f(1) = 1, f(2) = 2, f(3) = 3, f(4) = 4, f(5) = -10.\]

        \item Дан многочлен $f(x)$ степени не выше $n$.
        Докажите, что для любых различных вещественных чисел $x_0, x_1, \dotsc , x_n$ существуют и единственные $A_0, A_1, \dotsc , A_n$ такие, что

        \[\frac{f(x)}{(x - x_0)(x-x_1)\dotsc(x-x_n)} = \frac{A_0}{x-x_0} + \frac{A_1}{x-x_1} + \dotsc \frac{A_n}{x-x_n}\]

        \item Числа $p, q, r$ и $s$ различны.
        Докажите, что система линейных уравнений имеет единственное решение относительно переменных $x, y, z, u$ при любых значениях свободных членов $p_1, q_1, r_1, s_1.$

        \[
            \begin{cases}
                p^{3}x+p^{2}y+pz+u = p_1\\
                q^{3}x+q^{2}y+qz+u = q_1\\
                r^{3}x+r^{2}y+rz+u = r_1\\
                s^{3}x+s^{2}y+sz+u = s_1
            \end{cases}\,.
        \]

        \item Многочлен $f(x)$ принимает рациональные значения во всех рациональных точках.
        Докажите, что все коэффициенты $f(x)$ рациональны.

        \item Многочлен $f(x)$ степени не выше чем $n$ принимает в точках $0, 1, \dotsc, n$ значения $2^0, 2^1, \dotsc, 2^n$ соответственно.
        Найдите

        \begin{enumerate}
            \item $f(n+1);$
            \item $f(-1).$
        \end{enumerate}
        \item
        \begin{enumerate}
            \item Многочлен $f(x)$ степени не выше чем $n$ принимает целые значения в точках $0, 1, 2, \dotsc, n$.
            Докажите, что он принимает целые значения во всех целых точках.
            \item Многочлен $f(x)$ степени не выше чем $n$ принимает целые значения в точках $0, 1, 4, \dotsc , n^2$.
            Докажите, что он принимает целые значения во всех точных квадратах.
        \end{enumerate}
        \item Дан унитарный многочлен $P(x)$ степени $n$ и различные целые числа $x_0, x_1, \dotsc, x_n$.
        Докажите, что найдется такое $k$, что $|P(x_k)| > \dfrac{n!}{2^n}$

    \end{enumerate_boxed}

    \begin{figure}[h]
        \center{\includegraphics[width=0.875\linewidth]{meme}}\label{fig:meme}
    \end{figure}
\end{document}